\section{The Fermat character and the first variation}
\label{subsec:fermat-first-variation}

Throughout this section, characters refer to coordinate pullback.
Let $k=\overline{\mathbb F}_2$, let $r\geq3$ be odd, and put
\[
 H=z_1^r+\cdots+z_r^r,\qquad X_0=V(H)\subset\mathbb P_k^{r-1}.
\]
The group $D=\mu_r^r$ acts by $z_i\mapsto\lambda_i z_i$.
Its diagonal subgroup acts trivially on projective space.

\begin{proposition}[The top coherent character]
Put $E_0=H^{r-2}(X_0,\mathcal O_{X_0})$.
The character of $D$ on this line is
\[
 \chi_{\rm top}(\lambda_1,\ldots,\lambda_r)
       =(\lambda_1\cdots\lambda_r)^{-1}.
\]
This character descends to $D/\mu_r$.  Absolute $p$-power
Frobenius carries a $\chi$-eigenvector to a $\chi^p$-eigenvector.
In particular, for $p=2$ the absolute Frobenius on $E_0$ is zero.
\end{proposition}
\begin{proof}
Since $H$ is $D$-invariant, the hypersurface sequence
\[
 0\longrightarrow\mathcal O_{\mathbb P^{r-1}}(-r)
   \xrightarrow{\ H\ }\mathcal O_{\mathbb P^{r-1}}
   \longrightarrow\mathcal O_{X_0}\longrightarrow0
\]
is equivariant.  Its connecting map identifies $E_0$ with
$H^{r-1}(\mathbb P^{r-1},\mathcal O(-r))$, whose standard
\v Cech generator is $e=(z_1\cdots z_r)^{-1}$.
Pullback multiplies this generator by $(\prod_i\lambda_i)^{-1}$.
On a diagonal element $(\lambda,\ldots,\lambda)$ this scalar is
$\lambda^{-r}=1$, which proves descent.

Naturality of absolute Frobenius and its semilinearity give
$g^*F(e)=F(g^*e)=\chi(g)^pF(e)$.  A nonzero Frobenius image
in the same one-dimensional space would therefore require
$\chi^p=\chi$.  The displayed character has order $r$, whereas
for $p=2$ this requirement would make it trivial.  As $r\geq3$,
the Frobenius image is zero.
\end{proof}

Let now $F$ be any degree-$r$ form over $\mathbb F_2$, let
$c=[z_1\cdots z_r]F\in\mathbb F_2$, and consider the smooth open
part of the pencil
\[
 X_u=V(H+uF),\qquad
 U=\operatorname{Spec}\mathbb F_2[u,\Delta(u)^{-1}],\quad
 \Delta(0)=1.
\]
The earlier parity-counting pencil is a special case \citep{ParityCompanion}.
Here the existence of $U$ can be seen directly. Put $d_0=r(r-2)+1$.
Every degree-$d_0$ monomial is divisible by some $z_i^{r-1}$.
For each such monomial, choose one such index and multiply
$\partial_i(H+uF)$ by the complementary monomial. In the basis of
all degree-$d_0$ monomials, the resulting square matrix is $I+uB$,
because $r=1$ in characteristic two. Its determinant $\Delta$ has
constant term one. After inverting $\Delta$, every degree-$d_0$
monomial belongs to the partial-derivative ideal, so the derivatives
have no common projective zero. This proves smoothness over $U$.
The hypersurface sequence over $U$ and its Laurent monomial
calculation give a canonical frame $e$ of
$E=R^{r-2}\pi_*\mathcal O_{X_U}$.

\begin{proposition}[The precise first-order statement]
\label{prop:correct-first-variation}
The linearized coherent Frobenius is the bundle map
\[
 \varphi:F_U^*E\longrightarrow E,\qquad
 \varphi(1\otimes e)=uc\,e.
\]
Its derivative at $u=0$ in this frame and parameter is $c$;
the expression $uc$ is its first-order term.  Intrinsically its
first variation belongs to
\[
 T_0^*U\otimes\operatorname{Hom}_k(F_k^*E_0,E_0).
\]
It is not the differential of the absolute Frobenius morphism
of the base, and it is not an ordinary endomorphism of $E_0$.
\end{proposition}
\begin{proof}
The coherent Frobenius formula identifies its matrix coefficient
with $[z_1\cdots z_r](H+uF)=uc$.  This formula follows directly
by applying Frobenius to the hypersurface sequence and the displayed
\v Cech generator; in the smooth setting it is also Katz's
Algorithm~2.3.7.14, pp.~35--36, and Special Case~2.3.7.17,
p.~36 \citep{Katz1972}.  The coefficient of the product monomial in
$H$ is zero, so $\varphi_0=0$.  The first jet of a bundle map
vanishing at the base point is the displayed tensor, proving
the intrinsic assertion.

For clarity, under a frame change $e'=a(u)e$, with $a(0)\ne0$,
the coefficient changes to $a(u)^{p-1}h(u)$, where $p=2$ and
$h(u)=uc$.  Since $h(0)=0$, its derivative changes by
$a(0)^{p-1}$, exactly the change of basis in the indicated
Frobenius-twisted Hom space.  A change of parameter is accounted
for by the cotangent factor.  On the dual numbers
$k[\epsilon]/(\epsilon^2)$, the semilinear map is explicitly
\[
  F((a+b\epsilon)e)=a^2\epsilon c\,e.
\]
By contrast, the differential of absolute base Frobenius
$u\mapsto u^2$ is zero.  These are different operations.
\end{proof}

\paragraph{Equivariance and the limit of the character interpretation.}
For arbitrary $F$, the full diagonal group does not preserve
the generic fiber of the pencil: it sends the direction $F$ to
$g^*F=F(\lambda_1z_1,\ldots,\lambda_rz_r)$.
There is consequently no full-$D$ character decomposition on
the generic fiber defined by this action.  What is well-defined
at the Fermat point is the twisted first variation of
Proposition~\ref{prop:correct-first-variation}.
Its covariance is explicit:
\[
 c(g^*F)=(\lambda_1\cdots\lambda_r)c(F)
             =\chi_{\rm top}(g)^{-1}c(F).
\]
Indeed, among degree-$r$ monomials, the only one of full-$D$
weight $\prod_i\lambda_i$ is $z_1\cdots z_r$: every exponent
must be congruent to one modulo $r$, and their sum is $r$.
Equivalently, the coordinate change $X_{g^*F}\to X_F$ pulls
the canonical frame back to $\chi_{\rm top}(g)e$; compatibility
with the semilinear map gives
$h_{g^*F}=\chi_{\rm top}(g)^{1-p}h_F$.
For $p=2$, the Hom space
$\operatorname{Hom}(F_k^*E_0,E_0)$ has character
$\chi_{\rm top}^{-1}$, as required.

The incidence form has a useful smaller symmetry.
Write it as $F=z_0P$, with $P$ homogeneous of degree two in
each of its $M$ pairs of incidence variables, so $r=2M+1$.
Scale $z_0$ by $\lambda_0$ and both coordinates of pair $e$
by $\lambda_e$.  Then
\[
 g^*F=\rho(g)F,\qquad
 \rho(g)=\lambda_0\prod_{e=1}^M\lambda_e^2
         =\chi_{\rm top}(g)^{-1}.
\]
Together with $u\mapsto\rho(g)^{-1}u$, this defines an action
on the total pencil.  Its subgroup $\ker\rho$ acts fiberwise
and has trivial character on the coherent line.  Thus this
symmetry gives an equivariant family with a nontrivial action
on its parameter, not a full diagonal character action on
every fiber.  The statement concerns the usual incidence form;
the general first-variation proposition also applies to any exceptional cubic chosen in a counting construction.

In particular, the literal assertion that ``$uc$ is the derivative
on the $\chi_{\rm top}$-line'' needs two corrections: the derivative
is $c$, and its source is the Frobenius twist of that line.
Identifying this coherent first variation with a deformation of
a fixed rank-one Jacobi-sum motive would require an additional
motivic construction; the character calculation alone supplies
no such identification.

\subsection{The rank parameter depends on the input model}
\label{subsec:rank-model}

A rank bound does not specify how Frobenius is presented.
If an $R\times R$ Frobenius matrix is given modulo $p$, its
trace requires $R$ diagonal entries and $R-1$ field additions.
By contrast, a succinct geometric formula need not supply a
basis or the matrix.  The rank-one coherent pencil above is
an example of this distinction, but coherent rank is not the
rank of the full pure motive.  The following gives a separate
obstruction for genuine Artin motives, with its syntax stated
explicitly.

Here $\mathrm{NP}$ is nondeterministic polynomial time, and
$\mathrm{RP}$ is randomized polynomial time with one-sided error:
inputs outside the language are always rejected, while inputs in the
language are accepted with probability at least one half.

\begin{proposition}[A promised bounded-rank geometric model]
Fix a finite field $\mathbb F_q$ of characteristic $p$.
Suppose a deterministic algorithm, on polynomial presentations of finite
\'etale schemes $Y/\mathbb F_q$, correctly computes the Frobenius
trace modulo $p$ whenever the rank of $H^0(Y)$ is at most one,
and runs in time polynomial in the presentation size on that
promised class.  Then $\mathrm{NP}=\mathrm{RP}$.
The rank here is a promised upper bound, not an additional
input revealing the exact number of points.
\end{proposition}
\begin{proof}
For a Boolean formula $\Phi$, use Boolean input variables and
one uniquely determined Boolean wire variable for each gate.
Impose $x_i^2-x_i=0$ on inputs and the usual equations
$t-ab=0$, $t-(a+b-ab)=0$, and $t-(1-a)=0$ for AND, OR, and
NOT gates, followed by the equation that the output equals one.
This is a polynomial-size, degree-at-most-two presentation.
Eliminating wire variables identifies its coordinate algebra
with the quotient of
\[
 \mathbb F_q[x_1,\ldots,x_n]/(x_i^2-x_i)_i
       \simeq\prod_{a\in\{0,1\}^n}\mathbb F_q
\]
which retains exactly the satisfying factors.  Hence $Y_\Phi$
is a finite \'etale scheme consisting of $S_\Phi$ rational points.
Its full $H^0$ has rank $S_\Phi$ and Frobenius acts as the
identity, so its trace is $S_\Phi$.

On the promise $S_\Phi\in\{0,1\}$, the proposed algorithm
distinguishes the two cases.  Valiant--Vazirani isolation
therefore implies $\mathrm{NP}=\mathrm{RP}$ (Valiant \& Vazirani,
1986, Theorem~1.1 and Corollary~1.2, printed p.~87).
One may clock the algorithm at its polynomial bound on every
input.  An unsatisfiable input stays unsatisfiable under every
isolation restriction; a satisfiable input becomes uniquely
satisfiable with inverse-polynomial probability.  Arbitrary
answers outside the promise do not create a false positive
from an unsatisfiable formula.  Repetition proves the stated
one-sided-error consequence.
\end{proof}

\begin{proposition}[Known constant rank with succinct correspondences]
\label{prop:constant-artin-rank}
Consider the following larger input model for rational Artin
motives.  Finite \'etale schemes are given by polynomial
presentations, including fixed non-split extensions
$\operatorname{Spec}\mathbb F_{q^\ell}$ and explicitly presented
fiber products with these extensions.  Pullback and pushforward along explicitly
presented finite maps, rational linear combinations,
composition, direct sums, and projector summands may be
specified by an unexpanded expression.  Validity of the
projectors is a promise.  In this model a deterministic algorithm computing
Frobenius trace modulo $p$ in time polynomial in expression
size, full rank, and $\log q$ would imply
$\mathrm{NP}=\mathrm{RP}$, even when the exact rank is given
and is the constant one for odd $p$, or the constant two for
$p=2$.
\end{proposition}
\begin{proof}
We describe all realizations explicitly, so no conjectural
motivic comparison is involved.  For a finite \'etale scheme
$Y$, its rational Artin realization is the permutation
representation on $Y(\overline{\mathbb F}_q)$; rational
idempotent correspondences define its direct summands.
Let $\pi:Y_\Phi\to\operatorname{Spec}\mathbb F_q$ be the
scheme just constructed.  The composite
\[
 \epsilon=\pi_*\pi^*:\mathbf1\longrightarrow\mathbf1
\]
is multiplication by $S_\Phi$, because pullback sends $1$ to
the constant function $1$ on all points and pushforward sums
its values.  Its expression has polynomial size without
enumerating those points.

Choose $\ell=2$ if $p$ is odd and $\ell=3$ if $p=2$, and let
$r=\ell-1$.  The permutation Artin motive of
$\operatorname{Spec}\mathbb F_{q^\ell}$ has a rational
augmentation summand $V$, obtained by subtracting its constant
projector.  It has rank $r$.  Frobenius is an $\ell$-cycle on
the full permutation representation, of trace zero; its
constant summand has trace one.  Consequently $V$ has trace
$-1$.  Its projector is
$1-\ell^{-1}\pi_\ell^*\pi_{\ell *}$, a fixed-size expression.

Here the action of the unevaluated scalar $\epsilon$ on $V$
also has an explicit allowed expression.  Write
$L=\operatorname{Spec}\mathbb F_{q^\ell}$ and let
$\pi_L:Y_\Phi\times L\to L$ be the base-changed finite map.
The polynomial presentation of this fiber product has
polynomial size, and
\[
 \epsilon_L=\pi_{L*}\pi_L^*
       =S_\Phi\operatorname{id}_{h(L)}.
\]
If $e_V$ denotes the displayed augmentation projector, then
$e_V\epsilon_L e_V$ is precisely the expression denoted
$\epsilon\operatorname{id}_V$ below.  Thus no evaluation of
$S_\Phi$ and no unspecified operation of transferring an
unevaluated scalar between objects is required.

On the promise $S_\Phi\in\{0,1\}$, both $\epsilon$ and
$1-\epsilon$ are idempotents.  Define the genuine projector
motive
\[
 M_\Phi=
 (\mathbf1^{\oplus r},(1-\epsilon)\operatorname{id})
 \ \oplus\ (V,\epsilon\operatorname{id}).
\]
If $S_\Phi=0$, it is $\mathbf1^{\oplus r}$; if $S_\Phi=1$,
it is $V$.  Thus its exact rank is always the known constant
$r$, and its trace is
\[
 \operatorname{Tr}(\operatorname{Frob}_q\mid M_\Phi)
       =r-\ell S_\Phi.
\]
The two residues are distinct modulo $p$ because $p\nmid\ell$.
All displayed expressions have polynomial size.  The same
isolation argument proves the complexity implication.
\end{proof}

The last proposition is specific to the stated succinct
correspondence syntax.  If the zero-cycle $\pi_*[Y_\Phi]$ must
instead be expanded as an integer multiple of the point, its
coefficient is already $S_\Phi$: producing that input carries
the counting problem.  Likewise a given exact rank
$S_\Phi$ trivializes the first proposition's particular
zero-dimensional example.  If the permitted objects and spans
are all split finite \'etale schemes, every resulting Artin
representation is trivial, and its trace equals its given
rank.  The non-split constant block and its base-changed spans
are therefore substantive assumptions of
Proposition~\ref{prop:constant-artin-rank}.
These distinctions prevent an
unqualified answer to the polynomial-rank question.  They
give a conditional obstruction for a broad formula model,
while leaving more restrictive, explicitly specified geometric
models to be analyzed separately.

